\documentclass[UTF8,12pt]{ctexart}
\usepackage[paperwidth=240mm,paperheight=200mm,margin=14mm,top=8mm]{geometry}
\usepackage{amsmath,amssymb,mathtools}
\usepackage{xcolor}
\pagestyle{empty}
\setlength{\parindent}{0pt}
\setlength{\parskip}{0pt}
\newcommand{\qn}[1]{\textbf{\large #1}\ }
\xeCJKDeclareCharClass{CJK}{"2460 -> "24FF}
\begin{document}
\qn{1987.1}
设 $y=\ln(1+ax)$，其中 $a$ 为非零常数，则 $y'=$______，$y''=$______．

\vfill
\newpage
\qn{1987.2}
曲线 $y=\arctan x$ 在横坐标为 1 的点处的切线方程是______，法线方程是______．

\vfill
\newpage
\qn{1987.3}
积分中值定理的条件是______，结论是______．

\vfill
\newpage
\qn{1987.4}
$\lim\limits_{n\to\infty}\left(\dfrac{n-2}{n+1}\right)^n=$______．

\vfill
\newpage
\qn{1987.5}
$\int f'(x)\,\mathrm{d}x=$______，$\int_a^b f'(2x)\,\mathrm{d}x=$______．
求极限 $\lim\limits_{x\to0}\left(\dfrac{1}{x}-\dfrac{1}{\mathrm{e}^x-1}\right)$．
设 $\begin{cases} x=5(t-\sin t), \\ y=5(1-\cos t), \end{cases}$ 求 $\dfrac{\mathrm{d}y}{\mathrm{d}x}$，$\dfrac{\mathrm{d}^2y}{\mathrm{d}x^2}$．
计算定积分 $\int_0^1 x\arcsin x\,\mathrm{d}x$．
设 $D$ 是由曲线 $y=\sin x+1$ 与三条直线 $x=0$，$x=\pi$，$y=0$ 围成的曲边梯形，求 $D$ 绕 $Ox$ 轴旋转一周所生成的旋转体的体积．

\vfill
\newpage
\qn{1987.6}
若 $f(x)$ 在 $(a,b)$ 内可导，且导数 $f'(x)$ 恒大于零，则 $f(x)$ 在 $(a,b)$ 内单调增加．

\vfill
\newpage
\qn{1987.7}
若 $g(x)$ 在 $x=c$ 处二阶导数存在，且 $g'(c)=0$，$g''(c)<0$，则 $g(c)$ 为 $g(x)$ 的一个极大值．
计算不定积分 $\int\dfrac{\mathrm{d}x}{a^2\sin^2x+b^2\cos^2x}$，其中 $a,b$ 是不全为 0 的非负常数．

\vfill
\newpage
\qn{1987.8}
求微分方程 $x\dfrac{\mathrm{d}y}{\mathrm{d}x}=x-y$ 满足条件 $\left.y\right|_{x=\sqrt{2}}=0$ 的特解．

\vfill
\newpage
\qn{1987.9}
求微分方程 $y''+2y'+y=x\mathrm{e}^x$ 的通解．

\vfill
\newpage
\qn{1987.10}
$f(x)=|x\sin x|\,\mathrm{e}^{\cos x}$（$-\infty<x<+\infty$）是（　　）
\\
\noindent (A) 有界函数．　　(B) 单调函数．
\\
\noindent (C) 周期函数．　　(D) 偶函数．

\vfill
\newpage
\qn{1987.11}
函数 $f(x)=x\sin x$（　　）
\\
\noindent (A) 当 $x\to\infty$ 时为无穷大．　　(B) 在 $(-\infty,+\infty)$ 内有界．
\\
\noindent (C) 在 $(-\infty,+\infty)$ 内无界．　　(D) 当 $x\to\infty$ 时有有限极限．

\vfill
\newpage
\qn{1987.12}
设 $f(x)$ 在 $x=a$ 处可导，则 $\lim\limits_{x\to0}\dfrac{f(a+x)-f(a-x)}{x}$ 等于（　　）
\\
\noindent (A) $f'(a)$．　　(B) $2f'(a)$．
\\
\noindent (C) $0$．　　(D) $f'(2a)$．

\vfill
\newpage
\qn{1987.13}
设 $I=t\int_0^{\frac{s}{t}}f(tx)\,\mathrm{d}x$，其中 $f(x)$ 连续，$s>0$，$t>0$，则 $I$ 的值（　　）
\\
\noindent (A) 依赖于 $s,t$．　　(B) 依赖于 $s,t,x$．
\\
\noindent (C) 依赖于 $t,x$，不依赖于 $s$．　　(D) 依赖于 $s$，不依赖于 $t$．
在第一象限内求曲线 $y=-x^2+1$ 上的一点，使该点处的切线与所给曲线及两坐标轴所围成的图形面积为最小，并求此最小面积．

\vfill
\newpage
\end{document}